% Generated draft: requires mathematical and source review.
\documentclass[11pt]{article}
\usepackage[margin=25mm]{geometry}
\usepackage{fontspec}
\setmainfont{Latin Modern Roman}
\usepackage{amsmath,amssymb,mathtools,unicode-math}
\setmathfont{Latin Modern Math}
\usepackage{enumitem,needspace,xurl,hyperref}
\setlength{\parindent}{0pt}
\setlength{\parskip}{5pt}
\setlength{\emergencystretch}{3em}
\newcommand{\N}{\mathbb{N}}
\newcommand{\Z}{\mathbb{Z}}
\newcommand{\Q}{\mathbb{Q}}
\newcommand{\R}{\mathbb{R}}
\newcommand{\C}{\mathbb{C}}
\newcommand{\F}{\mathbb{F}}
\newcommand{\A}{\mathbb{A}}
\newcommand{\PP}{\mathbb{P}}
\newcommand{\E}{\mathbb{E}}
\newcommand{\eps}{\varepsilon}
\newcommand{\dd}{\,\mathrm{d}}
\begin{document}
\section*{UnsolvedMath ID: 30005704}
\textbf{Source problem code:} catalogue entry
\subsection{Definitions and mathematical statement}
A totally disconnected locally compact (TDLC) group is a Hausdorff locally compact topological group whose connected components are singletons. For a prime number \(p\), a finite elementary abelian \(p\)-subgroup of a group \(G\) is a finite subgroup isomorphic to \((\Z/p\Z)^r\) for some finite nonnegative integer \(r\). Including or excluding the trivial subgroup does not affect the finiteness question below. Define
\[
\mathcal E_p(G)=\{E\leq G:E\text{ is a finite elementary abelian }p\text{-group}\}.
\]
The group \(G\) acts on \(\mathcal E_p(G)\) by conjugation, and \(\mathcal E_p(G)/G\) denotes the set of its orbits. Thus the proposed conclusion is
\[
\bigl|\mathcal E_p(G)/G\bigr|<\infty.
\]
For comparison, an abstract group \(\Gamma\) has type \(\mathrm{FP}_\infty(\Z)\) if its trivial \(\Z[\Gamma]\)-module \(\Z\) admits a projective resolution whose modules are finitely generated in every degree. This definition alone does not specify the intended topological analogue for a TDLC group.

\textbf{Formulation problem.} Formulate an analogue for TDLC groups of the discrete finite-subgroup conjugacy theorem referred to in the supplied record. The proposed analogue is to give appropriate TDLC hypotheses corresponding to membership in Kropholler's hierarchy and to type \(\mathrm{FP}_\infty(\Z)\), and ask whether these imply that \(\mathcal E_p(G)/G\) is finite. A possible logical shape, rather than a fully specified conjecture, is
\[
\text{for every prime }p\text{ and every TDLC group }G,\qquad
\left.
\begin{gathered}
G\text{ satisfies the intended hierarchy hypothesis},\\
G\text{ satisfies the intended finiteness hypothesis}
\end{gathered}
\right\}
\Longrightarrow
\bigl|\mathcal E_p(G)/G\bigr|<\infty.
\]
The supplied record does not define either hypothesis, identify the precise version of Kropholler's hierarchy, or reproduce the discrete theorem being generalized. It also does not explicitly settle the quantification over \(p\); the displayed universal quantification describes the natural per-prime target and requires confirmation. Consequently this is an unresolved formulation request, not a precise assertion for a determined class of TDLC groups. No bound uniform in \(p\), and no finiteness assertion combining all primes, is specified.
\subsection{Short English statement}
The question seeks suitable hypotheses on a totally disconnected locally compact group under which, for a prime \(p\), there are only finitely many conjugacy classes of finite elementary abelian \(p\)-subgroups. The required hierarchy and homological finiteness hypotheses are unspecified, so a complete conjecture cannot yet be stated. The supplied record labels the problem open; that status has not been independently verified.
\subsection{Sources}
\begin{itemize}
\item Local UnsolvedMath registry, supplied complete record for permanent ID \(30005704\), source problem code \texttt{catalogue entry}, title ``TDLC Analogue of Finite-Subgroup Conjugacy,'' in the Oberwolfach Reports Open Problems collection. No individual proposer is identified for this target.
\item Supplied source citation: \emph{Mini-Workshop: Homological aspects of TDLC-groups} (2024), Oberwolfach Reports, pp.\ 3103--3138, DOI \texttt{10.4171/owr/2023/55}. Supplied link: \url{https://doi.org/10.4171/owr/2023/55}. The record separately gives proposed year 2023; this draft retains the supplied year information without resolving it.
\item Supplied literature source, ``Oberwolfach report'': \url{https://ems.press/content/serial-article-files/48166}.
\item Supplied literature source, ``Castellano--Weigel'': \url{https://doi.org/10.1016/j.jalgebra.2016.01.008}. The record supplies no title or specific result from this reference establishing the intended hypotheses.
\item Supplied literature assessment, checked on \texttt{2026-08-22}: the discrete theorem is reported as quoted, while a TDLC analogue is explicitly posed without a general formulation or proof. The supplied status is ``open.'' This is supplied dated curation derived from the OWR record, not an independently verified literature assessment.
\item The registry labels its statement \texttt{corrected\_verified} and reports that it was checked against the official Oberwolfach report. That supplied verification claim is retained as provenance; no source was visited for this draft, and the formulation gaps described here remain unresolved.
\end{itemize}
\end{document}
